\subsection{The exponent-eleven quotient}
\label{subsec:eleven}

Put $C_{11}:v^2+v=u^{11}$ and $A_{11}=\Jac(C_{11})$.
This is the classical factor $J_{1,1,9}$: the substitution
$(x,y)=(-v,-u)$ identifies its function field with
$y^{11}=x(1-x)$. The map
\[
f_{11}:X_{11}\longrightarrow C_{11},\qquad
[X:Y:Z]\longmapsto(-XY/Z^2,X^{11}/Z^{11})
\]
is nonconstant and extends over the smooth proper models, by the
same direct identity as in \eqref{eq:maps}.
\citet[\S4, table on p.~219]{GrossRohrlich1978} records the finiteness
of $A_{11}(\Q)$ proved by Faddeev. Finiteness of this Mordell--Weil
group by itself does not imply the required condition on $\Sha$.
We establish that condition using the central value and the published
Hilbert-modular theorem.

\begin{proposition}\label{prop:eleven}
The groups $A_{11}(\Q)$ and $\Sha(\Q,A_{11})$ are finite,
$A_{11}(\Q)\simeq\Z/11\Z$, and
\[
C_{11}(\Q)=\{\infty,(0,0),(0,-1)\}.
\]
The central value satisfies
\[
1.78<L(A_{11}/\Q,1)<1.80.
\tag{5.2}\label{eq:elevenvalue}
\]
\end{proposition}

\begin{proof}
We first fix the analytic normalization and certify nonvanishing.
\subsubsection{The analytic inputs}

Write $K=\Q(\zeta_{11})$, $F=K^+$, and
$\lambda=1-\zeta_{11}$. The Jacobi-sum description identifies
$L(A_{11}/\Q,s)$ with one Hecke $L$-function on $K$.
For a prime ideal $\mathfrak p\nmid11$, its value is
\[
j(\mathfrak p)=-\sum_{x\ne0,1}\chi_{\mathfrak p}(x)
                                      \chi_{\mathfrak p}(1-x),
\qquad |\sigma j(\mathfrak p)|=\sqrt{N\mathfrak p}
\]
for every complex embedding $\sigma$.
See \citet[\S1 and \S3]{GrossRohrlich1978} and
\citet[\S2.1--2.2, pp.~5--6]{ShuExpanded2021}.
Here and below the precise proposition numbering for Shu is that of
the expanded arXiv version; the root-number results are published in
\citet{Shu2021}.

The character has conductor $\lambda^2$ and the functional-equation
sign is $+1$. These inputs can be checked explicitly in
\citet[Proposition 3.12(2), p.~16; Theorem 3.13, pp.~17--18]{ShuExpanded2021}.
For $(r,s,t,\delta)=(1,1,9,1)$ put $q=(-2)^9$. Then
\[
v_{11}(q^{10}-1)=1,\qquad
 d=(q^{10}-1)/11\equiv
 9(2^{10}-1)/11\equiv1\pmod{11}.
\]
The conductor alternative is therefore $\lambda^2$, and the root-number
formula gives
\[
W=\left(\frac2{11}\right)
  \left[-\left(\frac{-2\cdot1\cdot1\cdot9\cdot d}{11}\right)\right]
  =(-1)(-1)=1.
\]
Since $|\operatorname{disc}K|=11^9$ and $N\lambda^2=11^2$, the rational
conductor is $11^{11}$. Thus, on putting
\[
B=\frac{\sqrt{11^{11}}}{(2\pi)^5},
\]
the completed function is
\[
\Lambda_{11}(s)=B^s\Gamma(s)^5L(A_{11}/\Q,s),\qquad
\Lambda_{11}(s)=\Lambda_{11}(2-s).
\tag{5.3}\label{eq:elevenfe}
\]
Entire continuation and the functional equation are the cited Hecke
inputs. In particular the numerical calculation below is not being
used to guess a conductor or a sign.

\subsubsection{A finite central-value certificate}

Let $G_5$ be the density on $(0,\infty)$ of the product of five
independent exponential random variables of mean one, and put
$Q_5(x)=\int_x^\infty G_5(u)\,du$. Its Mellin transform is
$\Gamma(s)^5$. Write $L(A_{11}/\Q,s)=\sum a_nn^{-s}$.
Mellin inversion and \eqref{eq:elevenfe}, split at $t=1$, give
\[
L(A_{11}/\Q,1)=2\sum_{n\geq1}\frac{a_n}{n}Q_5(n/B),
\qquad
Q_5(x)=\frac1{2\pi i}\int_{(c)}
             \frac{\Gamma(1+z)^5}{z}x^{-z}\,dz\quad(c>0).
\tag{5.4}\label{eq:elevenafe}
\]
For completeness the inverse Mellin series is
$\Theta(t)=\sum a_nG_5(nt/B)$, whose Mellin transform is
$\Lambda_{11}(s)$. The functional equation gives
$\Theta(t)=t^{-2}\Theta(1/t)$, so
$\Lambda_{11}(1)=2\int_1^\infty\Theta(t)\,dt$.
Substitution $u=nt/B$ yields the first formula in
\eqref{eq:elevenafe}. The positive-kernel bounds below justify absolute
convergence and the interchange of sums and integrals.

Let $b_K(n)$ count the integral ideals of $K$ of norm $n$.
The Hecke description and the triangle inequality imply
\[
|a_n|\leq \sqrt n\,b_K(n),\qquad
\sum_{n\geq1}b_K(n)n^{-\sigma}=\zeta_K(\sigma)\quad(\sigma>1).
\]
Markov's inequality for the product of the five exponential variables
uses its $k$-th moment $\Gamma(k+1)^5$ and gives
$Q_5(x)\leq\Gamma(k+1)^5x^{-k}$ for $k>0$.
Consequently, whenever $k+1/2>\sigma>1$,
\[
2\sum_{n>N}\frac{|a_n|}{n}Q_5(n/B)
 \leq 2\Gamma(k+1)^5B^k
       N^{\sigma-k-1/2}\zeta_K(\sigma).
\tag{5.5}\label{eq:eleventail}
\]
This follows term by term from
$n^{-k-1/2}\leq N^{\sigma-k-1/2}n^{-\sigma}$ for $n>N$.

The companion program computes every required $a_n$ exactly through
$N=100000$. Here are all the ingredients of that finite calculation.
For a rational prime $\ell\ne11$, put
$f=\operatorname{ord}_{11}(\ell)$. If $f$ is even, the Euler polynomial
is
\[
P_\ell(T)=(1+\ell^{f/2}T^f)^{10/f};
\]
this is \citet[Lemma 1.1]{GrossRohrlich1978}.
If $f=5$, it is $1-s_\ell T^5+\ell^5T^{10}$, where a direct count
of $C_{11}(\F_{\ell^5})$ gives
$s_\ell=(\ell^5+1-\#C_{11}(\F_{\ell^5}))/5$.
Only $\ell=3,5$ occur in this case below the cutoff. The resulting
polynomials are
\[
P_3(T)=1-24T^5+243T^{10},\qquad
P_5(T)=1-90T^5+3125T^{10}.
\]
If $f=1$, choose a primitive root of $\F_\ell^\times$ and count the
residue classes of
$\log(x)+\log(1-x)\pmod{11}$. If these counts are $c_0,\ldots,c_{10}$,
then $j=-\sum c_i\zeta_{11}^i$.
The Euler polynomial is $\prod_{a=1}^{10}(1-\sigma_a(j)T)$;
its required coefficients follow from Newton's identities using
integer arithmetic in $\Z[\zeta_{11}]$.
The factor at $11$ is one. When $\ell^f>N$, its entire contribution
through the cutoff is the constant term. Inversion of the computed
local polynomials and multiplicativity produce 1077 nonzero
coefficients, beginning with $a_1=1$, $a_{23}=21$, $a_{67}=-23$,
$a_{89}=21$, and $a_{243}=24$.

The finite sum in \eqref{eq:elevenafe} is evaluated by one contour
integral on $\Re z=2$:
\[
L_N=\frac2\pi\Re\int_0^\infty
 \frac{\Gamma(3+it)^5}{2+it}
 B^{2+it}\sum_{n\leq N}a_nn^{-3-it}\,dt.
\tag{5.6}\label{eq:elevencontour}
\]
The integral from zero to $20$ is enclosed with 96-bit Arb arithmetic
on twenty unit intervals. To bound the remainder, set
$S_N=\sum_{n\leq N}|a_n|n^{-3}$. The gamma integral gives
$|\Gamma(3+m+it)|\leq\Gamma(3+m)$; repeated use of its recurrence
therefore yields
$|\Gamma(3+it)|\leq\Gamma(3+m)|t|^{-m}$.
For every positive integer $m$, the contribution from $t\geq T$ is
at most
\[
\frac2\pi B^2S_N\frac{\Gamma(3+m)^5}{5m\,T^{5m}}.
\tag{5.7}\label{eq:elevenvertical}
\]
With $T=20,m=17$ this is less than $1.538\cdot10^{-24}$.
The needed Dedekind-zeta value is enclosed using
\[
\zeta_K(\sigma)=\zeta(\sigma)
 \prod_{\substack{\chi\bmod11\\\chi\ne1}}L(\sigma,\chi),\qquad
L(\sigma,\chi)=11^{-\sigma}\sum_{a=1}^{10}
                           \chi(a)\zeta(\sigma,a/11).
\]
The separate Riemann-zeta factor includes the prime at $11$; replacing
it by the imprimitive principal Dirichlet factor would omit that prime.
Arb certifies $\zeta_K(11/10)<3.007268$.
The numerical enclosures are
\[
\begin{split}
L_N&=1.789861079603+\varepsilon_N,
                    \qquad |\varepsilon_N|<1.54\cdot10^{-13},\\
|L(A_{11}/\Q,1)-L_N|&<0.004239358.
\end{split}
\]
The latter uses \eqref{eq:eleventail} with $k=4,\sigma=11/10$ and
\eqref{eq:elevenvertical}. These bounds prove \eqref{eq:elevenvalue}.

\subsubsection{The arithmetic conclusion}

Automorphic induction of the Jacobi character from $K$ to $F$ gives
the parallel-weight-two Hilbert form attached to $A_{11}/F$;
this is the realization in \citet[\S2.2]{ShuExpanded2021}.
Its central character is trivial: the inducing character restricts to the
quadratic character of $K/F$ on the real idele group, so the central
character of its automorphic induction is the square of that character
\citep[\S3.4, p.~15]{ShuExpanded2021}. One can also check this normalization
is that the real multiplication $F\subset\operatorname{End}_F^0(A_{11})$
is fixed by Rosati, so each two-dimensional coefficient component of
the polarized Tate representation has determinant the cyclotomic
character. The corresponding weight-two form therefore has trivial
finite central character. Each coefficient conjugate has the same
Hecke $L$-function after reindexing the prime ideals of $K$ by the
corresponding automorphism; Jacobi sums satisfy
$j(\tau\mathfrak p)=\tau j(\mathfrak p)$.
Thus \eqref{eq:elevenvalue} establishes nonvanishing for every
coefficient component of $A_{11}/F$.

Now $[F:\Q]=5$ is odd. Apply
\citet[Theorems B and B$'$, alternative (a), p.~593]{Nekovar2012}.
The vector-space Selmer groups vanish at every coefficient prime,
and the divisible-torsion Selmer groups vanish outside a finite set.
The Selmer exact sequence gives finite $A_{11}(F)$ and finite
$\Sha(F,A_{11})$, exactly as in the proof of
Proposition~\ref{prop:arithmetic}. Restriction to $F$ has a finite
kernel on $\Sha$, contained in
$H^1(\Gal(F/\Q),A_{11}(F))$, so the two groups over $\Q$ are finite.
This proves the new arithmetic requirement without a converse BSD
assumption. The torsion description of
\citet[Theorem 1.1]{GrossRohrlich1978} then gives
$A_{11}(\Q)\simeq\Z/11\Z$.

To recover the point set, put $P_0=(0,0)$ and $P_1=(0,-1)$.
The divisor of $v$ is $11P_0-11\infty$, and
$D=[P_0-\infty]$ generates $A_{11}(\Q)$.
The hyperelliptic involution gives $[P_1-\infty]=-D$.
At either $P_i$ the differentials
\[
\frac{du}{2v+1},\quad\frac{u\,du}{2v+1},\quad\ldots,
\quad\frac{u^4\,du}{2v+1}
\]
have vanishing orders $0,1,2,3,4$. Riemann--Roch therefore gives
$h^0(C_{11},dP_i)=1$ for $1\leq d\leq5$.
If a point $P$ represented $dD$, or $-dD$, with $2\leq d\leq5$,
then $P+(d-1)\infty\sim dP_i$ for one of the two $i$.
Its unique effective divisor is $dP_i$, a contradiction.
The only possible point classes are consequently $0,D,-D$, giving
the asserted three points. This completes the proof of
Proposition~\ref{prop:eleven}.
\end{proof}

\begin{remark}\label{rem:eleven-factor-only}
Only the factor $A_{11}$ is used. The whole Fermat Jacobian
$\Jac(X_{11})$ has positive Mordell--Weil rank: its rank is six by
\citet[Introduction and \S4]{GrossRohrlich1978}.
The target criterion in Stoll's theorem allows this distinction.
\end{remark}
